\documentclass{article}

\usepackage{fancyhdr}
\usepackage{amsmath}
\usepackage{amsthm}
\usepackage{amsfonts}
\usepackage{tikz}
\usepackage[plain]{algorithm}
\usepackage{algpseudocode}
\usepackage{amssymb}
\usepackage{enumitem}

\usetikzlibrary{automata,positioning}

\topmargin=-0.45in
\evensidemargin=0in
\oddsidemargin=0in
\textwidth=6.5in
\textheight=9.0in
\headsep=0.25in

\linespread{1.1}

\pagestyle{fancy}
\lhead{\hmwkAuthorName}
\chead{}
\rhead{\hmwkHeaderTitle}
\cfoot{\thepage}

\renewcommand\headrulewidth{0.4pt}
\renewcommand\footrulewidth{0pt}

\setlength\parindent{0pt}

\newcommand{\hmwkTitle}{Homework\ 7}
\newcommand{\hmwkDueDate}{Thursday, March 19, 2026,\ 11{:}59 pm (Thai time)}
\newcommand{\hmwkClass}{Data Structure\ and Abstractions}
\newcommand{\hmwkClassTime}{Section 2}
\newcommand{\hmwkClassInstructor}{Kanat Tangwongsan}
\newcommand{\hmwkAuthorName}{\textbf{Jiraroj Wiruchpongsanon (6781617)}}

\newcommand{\hmwkHeaderTitle}{Data Structure and Abstractions: Homework 7}

\title{
    \vspace{2in}
    \textmd{\textbf{\hmwkClass:\ \hmwkTitle}}\\
    \normalsize\vspace{0.1in}\small{Due\ on\ \hmwkDueDate}\\
    \vspace{0.1in}\large{\textit{\hmwkClassInstructor\ \hmwkClassTime}}
    \vspace{3in}
}

\author{\hmwkAuthorName}
\date{}

\renewcommand{\part}[1]{\textbf{\large Part \Alph{partCounter}}\stepcounter{partCounter}\\}

\newcounter{partCounter}
\newcommand{\solution}{\textbf{\large Solution}}

\newtheorem*{theorem}{Theorem}

\theoremstyle{remark}
\newtheorem*{remark}{Remark}

\begin{document}

\maketitle
\begin{center}
    \textit{Data Struct. and \{Abstractions, OOP\} (Term II/2025--26)}\\
    \textit{due: thu mar 19 @ 11:59pm}\\
    \textit{built on 2026/03/10 at 21:08:37}
\end{center}
\newpage

This assignment contains both programming and written questions. For the written portion, you will typeset your answer, produce a PDF file called \texttt{hw7.pdf}. No other format will be accepted. For the programming portion, you will follow the naming described in this handout. To hand in, zip your work as \texttt{a7.zip} and upload it to Canvas. Your zip will contain only the following files:

\begin{itemize}
    \item \texttt{hw7.pdf}
    \item \texttt{MultiwayMerge.java}
    \item \texttt{MakeTree.java}
    \item \texttt{Uproot.java}
    \item \texttt{Decor.java}
\end{itemize}

To get started, you’ll need the starter pack for this assignment. In it, you’ll find a class for binary tree nodes, which will be used in multiple problems across the assignment.

\subsection*{Collaboration}

We interpret collaboration very liberally. You may work with other students. However, each student must write up and hand in his or her assignment separately. Let us repeat: You need to write your own code. You must not look at or copy someone else’s code. You need to write up answers to written problems individually. The fact that you can recreate the solution from memory will be taken as proof that you actually understood it, and you may actually be interviewed about your answers.

\newpage

\section*{Task 1: Mathematical Truth (2 points)}

Remember binary trees from class? You will prove a proposition about them using a proof method of your choice:

\textbf{Proposition:} Every binary tree on $n$ nodes where each has either zero or two children has precisely $\frac{n+1}{2}$ leaves.

(Hints: By (strong) induction. It helps to recognize that a binary tree is either empty (null) or a node with two subtrees. This means that if a binary tree has $n$ nodes, then $n = 1 + n_L + n_R$, where $n_L$ and $n_R$ are the number of nodes in the left and right subtrees, respectively.)

\medskip
\solution

\begin{proof}

we prove this with strong induction, we claim that

\[P(n) \equiv \text{Every binary tree with n nodes, where each node has either 0 or 2 children has exactly } \frac{n+1}{2} \text{ leaves.}\]

for the \textbf{base case} where $n_0 = 0$ 

\[
\begin{aligned}
P(0) &= 0 = \left\lfloor{\frac{0+1}{2}}\right\rfloor \\
P(1) &= 1 = \left\lfloor{\frac{1+1}{2}}\right\rfloor = \left\lfloor{\frac{2}{2}}\right\rfloor \\
P(3) &= 2 = \left\lfloor{\frac{3+1}{2}}\right\rfloor = \left\lfloor{\frac{4}{2}}\right\rfloor
\end{aligned}
\]

as illustrated below, we'll see that the base case is satisfied --- regardless of what you consider as a bare minimum `binary tree` (n = 0 or n = 1)
 
\begin{center}
\begin{tikzpicture}[scale=1, every node/.style={circle, draw, minimum size=7mm, inner sep=0pt}]

% n = 0
\node[draw=none] at (0,2) {$n=0$};
\node[draw=none] at (0,0.8) {empty tree};

% n = 1
\node[draw=none] at (5,2) {$n=1$};
\node (a) at (5,0.8) {};
\node[draw=none] at (5,-0.2) {1 leaf};

% n = 2
\node[draw=none] at (10,2) {$n=2$};
\node (b1) at (10,1.1) {};
\node (b2) at (9.2,0.1) {};
\draw (b1) -- (b2);
\node[draw=none, align=center, text width=6cm] at (10,-2.0) 
{doesn't satisfy our predicate, as each of the binary tree's node could only have 0 or 2 child node};

% n = 3
\node[draw=none] at (15,2) {$n=3$};
\node (c1) at (15,1.2) {};
\node (c2) at (14.2,0.2) {};
\node (c3) at (15.8,0.2) {};
\draw (c1) -- (c2);
\draw (c1) -- (c3);
\node[draw=none] at (15,-0.6) {2 leaves};

\end{tikzpicture}
\end{center}

by observing the charasteristic of a binary tree (each node with 0 or 2 childrens), we can write the related of initial root note as 

\[n = 1 + n_l + n_r \tag{i}\]

where $n_l$ is the root node of all the nodes on the left, and $n_r$ is the root node of all the nodes on the right. we then assume $P(k)$ holds true for arbitrary $k \ge n_0$, and try to prove that $P(k) = \frac{n+1}{2}$

\vspace{1em}

as $n_l < n$ and $n_r < n$, $P(n_l)) \text{ and } P(n_r)$ is assumed to be true, and each side would have $\frac{n_l + 1}{2}$ and $\frac{n_r + 1}{2}$ accordingly. 

\vspace{1em}

so, according to our inductive hypothesis, binary tree with n nodes, where each node has either 0 or 2 children will has 

\[\frac{n_l + 1}{2} + \frac{n_r + 1}{2} = \frac{n_l + n_r + 2}{2}\]

and as we got the relation in (i), 

\[\frac{n_l + n_r + 2}{2} = \frac{n + 1}{2}\]

and the \textbf{inductive case} is satisfied, and we can conclude that every binary tree with n nodes, where each node has either 0 or 2 children, has exactly  $\frac{n+1}{2}$ leaves.

\end{proof}

\vspace{1em}
\hrulefill

\section*{Task 3: Let’s Grow a Tree (2 points)}

For this task, save your work in \texttt{MakeTree.java}

In this task, you will write a function to construct a ``balanced'' binary search tree from a list of keys---and analyze its running time.

\textbf{Subtask I:} You are to write a function

\begin{verbatim}
public static BinaryTreeNode buildBST(int[] keys)
\end{verbatim}

that takes an array of integer keys and returns a BST represented using the aforementioned format. If $n$ is the length of \texttt{keys}, your algorithm should take at most $O(n \log n)$ time and construct a BST that is no deeper than $1 + \log_2 n$ levels.

\textbf{Subtask II:} You will also analyze the running time of this algorithm and argue why the tree generated meets the depth requirement.

\medskip
\solution

\end{document}
